\documentclass[11pt,a4paper]{article}
\usepackage[margin=24mm,headheight=15pt]{geometry}
\usepackage{fontspec}
\setmainfont{texgyrepagella}[Extension=.otf,UprightFont=*-regular,BoldFont=*-bold,ItalicFont=*-italic,BoldItalicFont=*-bolditalic]
\setsansfont{texgyreheros}[Extension=.otf,UprightFont=*-regular,BoldFont=*-bold,ItalicFont=*-italic,BoldItalicFont=*-bolditalic]
\setmonofont{lmmono10-regular.otf}[Scale=0.86]
\usepackage{amsmath,amssymb,booktabs,longtable,tabularx,array,graphicx}
\usepackage{pdflscape,caption,fancyhdr,xcolor,enumitem,float}
\usepackage[numbers,sort&compress]{natbib}
\usepackage{hyperref,xurl}
\hypersetup{colorlinks=true,linkcolor=blue!45!black,citecolor=blue!45!black,urlcolor=blue!45!black,
pdftitle={SRG Simulation Reconstruction — Volume III: Band-Wide and Ordered Local Fusion},
pdfauthor={Hilmir Frímann Halldórsson; scientific lead GPT; manuscript preparation Codex},
pdfsubject={Candidate v0.1 for internal review; not publication-approved}}
\pagestyle{fancy}\fancyhf{}
\fancyhead[L]{\small SRG reconstruction: Tower II and Universe}
\fancyhead[R]{\small Candidate v0.1}
\fancyfoot[C]{\thepage}
\setlength{\parindent}{0pt}\setlength{\parskip}{6pt}
\setlength{\emergencystretch}{2em}
\renewcommand{\arraystretch}{1.18}
\captionsetup{font=small,labelfont=bf}
\setlist{nosep,leftmargin=*}
\newcolumntype{P}[1]{>{\raggedright\arraybackslash}p{#1}}
\newcommand{\ev}[1]{\textsf{\small #1}}
\newcommand{\src}[2]{\textsf{#1:#2}}
\newcommand{\R}{\mathbb R}
\newcommand{\sgn}{\operatorname{sgn}}
\newcommand{\clip}{\operatorname{clip}}
\newcommand{\mean}{\operatorname{mean}}
\newcommand{\sd}{\operatorname{std}}
\newcommand{\norm}[1]{\left\lVert #1\right\rVert}
\newcommand{\plate}[3]{\par\begin{figure}[!tbp]\centering
\includegraphics[width=\textwidth,height=.77\textheight,keepaspectratio]{#1}
\caption{#2}\label{#3}\end{figure}\par}
\newcommand{\wideplate}[3]{\clearpage\begin{landscape}\begin{figure}[!ht]\centering
\includegraphics[width=\linewidth,height=140mm,keepaspectratio]{#1}
\caption{#2}\label{#3}\end{figure}\end{landscape}\clearpage}
\begin{document}
\hypersetup{pageanchor=false}
\begin{titlepage}
\vspace*{14mm}
{\sffamily\large SRG SIMULATION RECONSTRUCTION\par}
\vspace{7mm}{\sffamily\Huge Volume III\par}
\vspace{8mm}{\sffamily\LARGE Band-Wide and Ordered Local Fusion\\[4pt]in Tower II and Universe\par}
\vspace{12mm}{\large Hilmir Frímann Halldórsson\par}
Author and research director\\
Scientific lead and reviewer: GPT\\
Research engineering and manuscript preparation: Codex
\vspace{9mm}

\textbf{Candidate v0.1 · 25 October 2026}\\
For GPT and Hilmir's scientific review. Volume number provisional.\\
This candidate has not been accepted for publication.

\vspace{8mm}\textbf{Abstract}\par
Two historical SRG programs move indexed points through local geometry changes, real-amplitude exchange, prescribed time modulation, center-directed displacement, signed compression and positional noise. Tower II calculates one weighted center for each permanently assigned band. Universe, when its horizon is finite, calculates a different local center at each ascending point ID and reads coordinates already changed earlier in that pass. This reconstruction makes those dependencies explicit and connects them to accepted records. Tower II's local graph separates while its band-center dependency union remains weakly connected throughout 600 inspected incoming boundaries. Universe's accepted 200-pass replay retains ten spatial-observer memberships from pass 16 onward, yet its selected edges and individual positions continue changing. Captured operands identify actual earlier writes in later centers; they do not measure a reversed-order alternative. Historical summaries, reproduced coordinates and reviewed mathematical identities are distinguished throughout. Results and figures are reused from three accepted internal reviews. Manuscript preparation introduces no simulation, numerical scientific rerun or new physical assumption.

\vfill\small
Historical source headers credit \emph{Claude (collaborating with pzychozen)} and state 14 March 2026. This attribution is preserved; it is not a verified execution chronology. Internal GPT acceptance of supporting studies is distinct from external peer review.
\end{titlepage}
\hypersetup{pageanchor=true}
\tableofcontents\clearpage
\section{What exists in these simulated worlds}\label{sec:orientation}
A point has a permanent row number, three coordinates and one real number called its resonance. Each update changes coordinates and resonance. The row number stays with the point even when plotted markers overlap. A group can mean points assigned the same frequency, points that started near one another, points connected by current neighbor selections, or points joined by a separate clustering measurement. Those definitions answer different questions.

The positive question of this volume is how geometry and amplitude evolution are connected by the centers that the programs actually calculate. In Tower II, all members of an assigned band contribute to a common center even when they occupy separate local neighborhoods. In finite-horizon Universe, a point uses only its selected neighbors; its center reads positions available at the instant its row is updated. Earlier rows may already have moved.

Figure~\ref{fig:towerworld} distinguishes visible concentration from retained identity. The same 80 Tower II rows are drawn at every boundary. Late marker overlap is a limitation of display scale, not a deletion or merger. Universe's world view appears in Figure~\ref{fig:uworld}; its populations remain far apart while smaller pieces gather within them.

\textbf{Conceptual route.} Read this section, the captions, Sections~\ref{sec:towerobs}, \ref{sec:original}, \ref{sec:umotion} and \ref{sec:comparison}. Motion episodes explain what the records establish before broader interpretation.

\textbf{Exact route.} Sections~\ref{sec:common}, \ref{sec:tower}, \ref{sec:universe} and \ref{sec:observers} specify snapshot timing, branches and equations. Section~\ref{sec:fidelity} separates recording fidelity from historical agreement. Appendix~\ref{app:evidence} and the claim matrix identify source files, retained tables and archived operands.

\begin{table}[H]\centering\small
\begin{tabularx}{\textwidth}{p{30mm}X}
\toprule Term & Meaning used here\\\midrule
Point / ID & Persistent array row; \(X_i\in\R^3\), \(R_i\in\R\).\\
Band & Tower II's fixed assignment \(b_i\), with supplied frequency \(\omega_i\).\\
Population & Universe's supplied 20-ID block P0--P5, used for description, not as a coupling label.\\
Selected neighbor & A row selected from incoming distances, at most five per point.\\
Weak / strong component & Connected ignoring edge direction / mutually reachable through directed selections.\\
Spatial-observer group & Universe's average-linkage partition at cut distance 2, measured after an update.\\
Fusion & Source name for center-directed position displacement with a signed amplitude coefficient.\\
Velocity \(V\) & Last completed displacement; never consumed as momentum by the next update.\\
Boundary / pass & Boundary \(s\) is the state after \(s\) updates; pass \(p\) takes boundary \(p-1\) to \(p\).\\
\bottomrule\end{tabularx}
\caption{Definitions separating the different meanings of motion and grouping.}
\end{table}

\wideplate{evidence/S6/source_package/figures/F01_fixed_world.png}
{\ev{REPRODUCED / DERIVED}. Tower II, seed 2038, 80 fixed IDs colored by assigned band. The accepted plot shows boundaries 0, 20, 100, 200, 400 and 600 on equal axes with fixed limits \(\pm5.5\) and camera. Boundaries 0--200 are a later reproduction under the historical configuration; 400 and 600 belong to the existing R1 extension. No original 600-pass coordinates were recovered. Late overlap hides small separations and motion. S6 F01, unchanged; retained R1 coordinates and S6 geometry~\citep{s6study,s6review}.}
{fig:towerworld}

\section{Evidence, provenance and scope}\label{sec:evidence}
Six labels apply to prose, equations, tables and captions:
\begin{description}[style=nextline,leftmargin=0pt,labelindent=0pt]
\item[\ev{ORIGINAL}] Unchanged historical source, numerical file or narrative.
\item[\ev{ATTRIBUTED}] Configuration, motivation, date or interpretation stated by a source or narrative.
\item[\ev{OBSERVED IN HISTORICAL DATA}] Values or relationships in retained original numerical outputs.
\item[\ev{REPRODUCED}] States and outputs from a later source-faithful execution with its recorded environment.
\item[\ev{DERIVED}] Reviewed source transcription, identity or passive description of specified records. Its input evidence class remains explicit.
\item[\ev{UNRESOLVED}] A question the source and finite evidence do not settle.
\end{description}
These labels do not imply that a reproduction replaces an original. Source determines operations; original tables determine what numbers survive; reproduced records reveal the later realization; accepted reviews qualify interpretation.

The three accepted archives are S6 (Tower II construction and fixed-record analysis), S7C (Universe construction and original aggregates), and S7R (Universe finite replay). They contain submitted material unchanged together with GPT's independent review and additional findings~\citep{s6review,s7creview,s7rreview}. References are self-contained research sources. Acceptance is internal scientific review, not external peer review. This manuscript is a new candidate.

\begin{table}[H]\centering\small
\begin{tabularx}{\textwidth}{p{28mm}X X}
\toprule & Tower II & Universe\\\midrule
Original program & 842 lines; 34,980 bytes & 778 lines; 30,165 bytes\\
Original data & Seven breathing-setting summary rows & Five CSVs, 233 rows; main table \(200\times15\)\\
Selected setting & \(N=80\), four bands, seed 2038, width 2, \(\epsilon=1\), no horizon & \(N=120\), seed 1618, six populations, \(h=6\), assigned frequencies, breathing\\
Historical duration & 200 passes; no individual coordinates & 200 passes, rows 1--200; no individual coordinates\\
Detailed record & Existing R1 boundaries 0--600; extension after 200 & Reference and capture, each boundaries 0--200\\
Current drafting & No new updates or scientific reruns & Same zero-execution boundary\\
\bottomrule\end{tabularx}
\caption{Evidence windows. Universe's twins validate recording of one setting; they are neither independent trials nor a single 400-pass trajectory.}
\end{table}

Tower II's record was selected through the retained R1 index, not a search for a preferred shape. Its 601 boundaries contain 48,080 point-state rows. The comparison window ends at 200; the same engine continues in the already existing record to 600. SIM06 only inspected that record. Its graph analysis covers incoming boundaries 0--599 supplying passes 1--600.

Universe's historical main table contains no boundary-zero observer row. The accepted replay executed exactly two initializations and 400 total updates: one 200-pass reference and one matching 200-pass capture. Each contains 24,120 point-state rows. Captured incoming graphs cover 0--199. Neither boundary 200's next graph nor pass 201 was evaluated~\citep{s7rprovenance,s7rreview}.

Headers and narratives state 14 March 2026 and credit Claude collaborating with pzychozen~\citep{towernote,universenote}. Names such as Part II and Infinity convey intended relationships; they do not certify invention order, runtime dates or identity of engines. Study numbers record investigation order. No wider archive chronology is manufactured here.

\section{Shared state, clock and stage notation}\label{sec:common}
Let \(X=X_{p-1}\), \(R=R_{p-1}\) denote incoming arrays. Both engines increment an integer counter once per direct update call and set \(t=\operatorname{float}(p)\). There is no separate \(dt\), animation callback clock, wrapping or boundary wall. Initialization ranges are sampling ranges. Anchors are \src{T}{129--216} and \src{U}{155--203}; T and U identify unchanged originals~\citep{towersource,universesource}.

We use \(Y\) for post-local positions, \(B\) for amplitudes after exchange and model-specific additions, \(Q\) for post-\(H\), pre-compression amplitudes, and \(Z\) for post-fusion, pre-noise positions:
\begin{equation}\label{eq:stages}
\begin{gathered}
(X,R)\longrightarrow(d,\mathcal N)\longrightarrow Y,\qquad
R\longrightarrow B\longrightarrow Q,\\
(Y,Q,\mathcal N)\longrightarrow Z,\qquad
R_p=\operatorname{RPCO}_{L_p}(Q),\\
X_p=Z+\xi_p,\qquad V_p=X_p-X.
\end{gathered}
\end{equation}
This is conceptual ordering, not spatial geometry. Tower II obtains its bound before the graph; Universe obtains its time-only bound after fusion. Both fuse using \(Q\), not final \(R_p\).

For selected neighbors \(\mathcal N_i\),
\begin{align}
\ell_i&=\begin{cases}
|\mathcal N_i|^{-1}\displaystyle\sum_{j\in\mathcal N_i}(X_j-X_i),&
\mathcal N_i\ne\varnothing,\\0,&\mathcal N_i=\varnothing,
\end{cases}\label{eq:lap}\\
Y_i&=X_i+0.05\ell_i+0.15\tanh(\ell_i).\label{eq:local}
\end{align}
Tanh is componentwise. Every \(\ell_i\) is formed from incoming \(X\) before the simultaneous array addition. No amplitude or inverse-distance weights enter this position mean.

Both sources use
\begin{align}
H(t)&=2e^{-0.5t}\cos t+e^{-0.3t}\sin t+
1.5e^{-0.7t}+e^{-0.4t},\label{eq:H}\\
H_{\max}&=\max_{t\in\mathcal T}|H(t)|+10^{-12},\quad
\mathcal T=\operatorname{linspace}(0,100,50000),\label{eq:Hmax}\\
Q_i&=B_i[1+0.05H(p)/H_{\max}].\label{eq:Q}
\end{align}
This normalization is a sampled maximum; its retained value is \(4.500000000001\), not newly evaluated here. The source calls the stage memory, but it evaluates prescribed integer time, not a history convolution. \(H\)'s sine/cosine use argument \(t\); the label-specific sine uses \(2\pi\omega_i p\). These angular conventions remain distinct.

The signed compressor and positional noise are
\begin{align}
\operatorname{RPCO}_L(q)&=\sgn(q)\exp\{\clip[\log(|q|+10^{-9}),-L,L]\},\label{eq:rpco}\\
\xi_{p,id}&\sim\mathcal N(0,\sigma_p^2),\qquad
\sigma_p=\frac{0.01}{1+0.01p},\quad d=1,2,3.\label{eq:noise}
\end{align}
The distribution describes the historical call, not draws performed here. Declining positional noise is not velocity damping. Continuation requires \(X,R,p\), fixed labels/configuration and RNG state; diagnostic \(V\) and histories do not steer the next pass. No implemented energy or action functional governs these stages.

\section{Tower II: a complete pass on paper}\label{sec:tower}
\subsection{Initialization and selection}
The constructor \src{T}{62--109} optionally seeds NumPy's global legacy stream, draws uniform \(X\), then uniform \(R\), and sets \(V=0,p=0\). Defaults are coordinate range \([-5,5]\), amplitude range \([0,0.3]\). For \(n_b\) bands it takes the integer list from \(-\lfloor n_b/2\rfloor\) to \(+\lfloor n_b/2\rfloor\), truncated to \(n_b\) entries, and assigns frequencies \(0.244\phi^n\), where \(\phi=(1+\sqrt5)/2\). It shuffles a repeated \(i\bmod n_b\) allocation once. Selected \(N=80,n_b=4\) gives 20 IDs per band and exponents \(-2,-1,0,1\). Labels never evolve.

Before distances, breathing sets
\begin{equation}\label{eq:LT}
L_p=9\left(1+\epsilon\,N^{-1}\sum_i|R_i|\right),\qquad
\epsilon=1\text{ in the selected setting}.
\end{equation}
It reads incoming \(R_{p-1}\). Off means \(L=9\); enabled values are appended to history. \(L\) is overwritten, not integrated as an independent oscillator.

Distances are \(d_{ij}=\norm{X_j-X_i}\). With finite horizon, candidates satisfy strict \(d_{ij}<h\), self is explicitly removed, and nearest five or fewer are selected. Without a horizon the source uses \texttt{argsort(dists[i])[1:6]} (\src{T}{111--127}), assuming the first sorted entry is self rather than explicitly excluding it. The accepted selected record has no self-selections or fifth/sixth cutoff ties; this does not repair the general coincidence caveat. Later loops recover selected IDs in ascending ID order with \texttt{where(adj[i]>0)}. Equations~\eqref{eq:lap}--\eqref{eq:local} give \(Y\).

\subsection{Exchange, imposed sine and additive terms}
Initialize \(B=R\). For nonempty rows,
\begin{equation}\label{eq:TW}
u_{ij}=\frac{\exp[-c_w\{\log(\omega_j/\omega_i)\}^2]}{d_{ij}+10^{-8}},
\qquad S_i=\sum_{j\in\mathcal N_i}u_{ij}.
\end{equation}
If \(S_i<10^{-10}\), skip the entire amplitude row; equality is active. Otherwise \(W_{ij}=u_{ij}/S_i\). Selected width \(c_w=2\) differs from default 3. Decreasing width broadens the raw log-frequency preference.

The first row operations are
\begin{align}
A_i&=R_i+0.15\left(\sum_jW_{ij}R_j-R_i\right),\\
B_i^{(0)}&=A_i[1+0.01\sin(2\pi\omega_i p)].\label{eq:Tsine}
\end{align}
All right-hand sides read incoming \(R\). Additions then occur \emph{after} the sine (\src{T}{169--188}). For each selected \(j\) with \(|\omega_i-\omega_j|>0.01\), visit \(q=-2,-1,1,2\) and add
\begin{equation}\label{eq:phi}
M^\phi_{ijq}=\eta e^{-5\{\log(\omega_i/\omega_j)-\log(\phi^q)\}^2}R_j .
\end{equation}
For that same \(j\), set \(f_+=\omega_i+\omega_j\), \(f_-=|\omega_i-\omega_j|\); for every catalogue target \(f\), in catalogue order, add
\begin{align}
M^+_{ijf}&=\tfrac12\eta R_j\,\mathbf1_{\{|f_+-f|/f<0.2\}},\\
M^-_{ijf}&=\tfrac12\eta R_j\,\mathbf1_{\{f_->0.01\}}
                             \mathbf1_{\{|f_--f|/f<0.2\}}.\label{eq:sumdiff}
\end{align}
Thus active \(B_i=B_i^{(0)}+\sum M^\phi+\sum(M^++M^-)\), with all additions under the outer unequal-frequency condition. Selected \(\eta=0.001\) is the default. Every qualifying target accumulates; matches are not exclusive and donors are not debited. Empty or guard-skipped rows retain \(R_i\), skipping sine and additions. The array is installed after all rows.

The equations retain mathematical terms; exact floating evaluation keeps the source's loop and reduction order. At frozen geometry, frequencies and pass, these terms are linear in incoming \(R\). The historical label nonlinear mixing does not override that conditional fact. The full moving system is nonlinear through geometry, selection, tanh, magnitude weights and compression. Additions do not create evolving frequency labels or establish a conserved energy transfer~\citep{s6review,s6study}.

\subsection{Band centers and completed observers}
After Equation~\eqref{eq:Q}, each fixed band \(G_b\) with at least two members computes one center:
\begin{align}
c_b&=\frac{\sum_{j\in G_b}(|Q_j|+10^{-8})Y_j}
          {\sum_{j\in G_b}(|Q_j|+10^{-8})},\label{eq:Tcenter}\\
Z_i&=Y_i+0.02Q_i(c_b-Y_i),\qquad i\in G_b.\label{eq:Tfusion}
\end{align}
Singletons skip this stage. A center includes self and every band member regardless of local components or horizon. It is computed once before the inner writes. Earlier bands cannot change this band's operands because the bands are disjoint. It is neither one global center nor a sequentially recomputed local center.

Magnitude \(Q\) determines positive center weights; signed \(Q_i\) determines displacement direction. Positive \(Q_i\) points toward the center, but convex interpolation also requires \(0\le0.02Q_i\le1\). Negative values move away and sufficiently large positives overshoot. The retained selected record lies in the positive small-coefficient regime; the implementation allows others.

Compression uses \(L_p\), noise is added, \(V=X_p-X\) is stored and completed arrays are measured (\src{T}{209--250}). Band summaries include mean/max/std of \(|R|\), count, frequency/index, coherence and mean GRF. Coherence is
\begin{equation}\label{eq:lock}
C_b=1-\frac{\sd_{i\in G_b}|R_{p,i}|}
              {\mean_{i\in G_b}|R_{p,i}|+10^{-10}}.
\end{equation}
First \(C_b>0.7\) with mean \(|R|>0.01\) is retained as lock time. It is an amplitude threshold, not a positional stop. Global metrics include mean/max \(|R|\) and count of band means above 0.01. R1's extra squared-amplitude sum is diagnostic, not an established energy.

\subsection{Other drivers in the complete source}
The selected breathing driver (\src{T}{376--472}) uses seven \(\epsilon\) values, 200 passes each, and writes the original seven-row table. None is run here. Other sections matter because their names can be mistaken for recurrence mechanisms.

The master-equation driver (\src{T}{262--368}) evolves first, then calculates
\[
\gamma_{\rm pred}(\omega)=
\frac{\zeta(2\log_\phi(\omega/0.244)+3)}{\pi e\phi}
\]
when the zeta argument exceeds 1.01, and compares it with lock times. It is not an update coefficient. The frequency-generation driver (\src{T}{480--588}) analyzes amplitude spectra under different addition strengths; assigned labels stay fixed. Its line-552 modulo test is not a general correct fold onto a real sampled spectrum's Nyquist interval. This is a retained interpretive limit, not a new spectral audit.

The tunneling driver (\src{T}{596--744}) changes bridge populations and bands, settles 50 passes, triples one population's amplitudes and records 100 more. Band centers remain available across local separation and no unperturbed paired response is supplied. The label does not establish isolated signal transmission. The phase-diagram driver (\src{T}{752--821}) summarizes band counts and widths. The main block invokes all five; the original has no animation/rendering routine. These complete-source distinctions supply context without new outcome analysis.

\section{Tower II: what the saved motion establishes}\label{sec:towerobs}
\subsection{Gathering across fixed bands}
The accepted R1 local graph starts as one weak/strong component and first weakly splits at boundary 10. At boundary 20 its six weak components have 14,6,10,8,10,32 members and each contains all four bands. Spatial gathering does not immediately form four band objects.

IDs 4,8,9 were chosen as the first three band-0 IDs, not by searching for behavior. At 20, ID 4 is in an eight-point local component; 8 and 9 are in the 32-point component. All three still enter one band-0 center. At 200, band 0's induced local graph has nine weak pieces. The whole graph has 400 directed selections: 138 same-band and 262 cross-band. That graph supplies existing R1 extension pass 201, not an additional historical update~\citep{s6study,s6review}.

\begin{table}[H]\centering\small
\begin{tabular}{rrrrr}\toprule
Boundary & Whole \(R_g\) & Band-0 \(R_g\) & Mean \(|R|\) & Mean completed step\\\midrule
0 & 5.2615661 & 5.1145139 & .151692819 & 0 (initial \(V\))\\
20 & 4.2927015 & 4.2177045 & .181526538 & .05595313\\
200 & 1.3528301 & 1.2898877 & .548983269 & .01440317\\
600 (extension) & .05499544 & .01749509 & 2.182252330 & .00252415\\
\bottomrule\end{tabular}
\caption{\ev{REPRODUCED / DERIVED}. Rounded accepted summaries from S6 \protect\path{independent/FINDINGS.json}; endpoint reduction is not monotonicity.}\label{tab:Towergeometry}
\end{table}

Whole radius falls 74.29\% from 0 to 200 and 98.95\% from 0 to 600. Over 200--600, band 0's radius decreases by a factor about 73.73 while its arithmetic center translates only about 0.096678 between endpoints. Spread and center translation are different measurements. Mean completed step remains 0.00252415 at 600, maximum 0.00468384.

\wideplate{evidence/S6/source_package/figures/F02_band_and_point_detail.png}
{\ev{REPRODUCED / DERIVED}. Top: band 0 at 20,200,600 centered and divided by its current arithmetic \(R_g\), fixed normalized limits, no fitted rotation; printed radii disclose actual contraction. Bottom: IDs 4,8,9 over 0--20,20--200,200--600, equal aspect but separately disclosed world spans. Hollow markers start, filled markers end. Lines join saved samples only. Last interval wholly R1 extension. S6 F02, unchanged~\citep{s6study}.}
{fig:towerdetail}

\subsection{Local separation with persistent band dependencies}
The late example is incoming boundary \textbf{599}, supplying pass \textbf{600}:
\begin{table}[H]\centering\small
\begin{tabularx}{\textwidth}{r rrrr X}\toprule
Size & Band 0 & Band 1 & Band 2 & Band 3 & Band-0 IDs\\\midrule
41 & 1 & 20 & 0 & 20 & 62\\
22 & 2 & 0 & 20 & 0 & 25,74\\
17 & 17 & 0 & 0 & 0 & 4,8,9,11,12,15,17,18,20,31,34,45,55,57,66,67,70\\
\bottomrule\end{tabularx}
\caption{\ev{REPRODUCED / DERIVED}. Exact accepted boundary-599 composition. Band 0 spans all three local components. Complete memberships: S6 \protect\path{independent/FINDINGS.json}, \protect\path{snapshot_graphs/599}.}\label{tab:bridge}
\end{table}
GPT's additional accepted analysis takes the union of local directed selections and all ordered distinct same-band pairs. That analysis-only dependency graph is weakly connected at \emph{all 600} inspected incoming boundaries. Its strong-component count ranges from one to three; at 599 it is two~\citep{s6review}.

The union describes dependencies already in Equations~\eqref{eq:TW}--\eqref{eq:Tfusion}; it adds no evolution edges. It is not a full state-variable Jacobian, a measurement of every edge's nonzero influence, physical propagation speed or universal strong connectivity. Local component separation alone cannot establish independence of the full Tower II update.

\subsection{Within-band spread and center separation}
With arithmetic whole center \(\mu\), arithmetic band centers \(\mu_b\), and counts \(n_b\),
\begin{align}
R_g^2&=N^{-1}\sum_i\norm{X_i-\mu}^2=W+B_{\rm ctr},\label{eq:variance}\\
W&=\sum_b\frac{n_b}{N}\frac1{n_b}\sum_{i\in G_b}\norm{X_i-\mu_b}^2,\nonumber\\
B_{\rm ctr}&=\sum_b\frac{n_b}{N}\norm{\mu_b-\mu}^2.\nonumber
\end{align}
\(B_{\rm ctr}\) is descriptive variance, distinct from stage amplitude \(B_i\). Write \(X_i-\mu=(X_i-\mu_b)+(\mu_b-\mu)\), expand the square and sum within each band. The cross term vanishes since \(\sum_{i\in G_b}(X_i-\mu_b)=0\). The weights preserve point counts. These arithmetic centers differ from the \(Q\)-weighted fusion centers.

The between-center fraction increases from 3.78\% initially to 6.66\% at 200 and 96.65\% at 600, while \emph{both} absolute variance terms decrease. Most remaining late spread lies between compact bands; the fraction does not show centers separating farther. Figure~\ref{fig:towerdiag} plots radii whose squares add.

Lock times are band 2 at pass 3, band 3 at 4, band 1 at 5, band 0 at 7, all before the first weak split at 10. The amplitude-threshold chronology is neither a geometric formation time nor the end of motion.

\plate{evidence/S6/source_package/figures/F03_geometry_and_diagnostics.png}
{\ev{REPRODUCED / DERIVED}. S6 F03, unchanged. Shading after 200 marks R1 extension. Upper left: \(R_g,\sqrt W,\sqrt{B_{\rm ctr}}\), whose squares add. Upper right: separate axes for mean \(|R|\) and \(L_p\), which reads state \(p-1\). Middle: incoming graphs 0--599. Bottom: arithmetic-center translations and completed steps. Continued motion persists below world-marker scale~\citep{s6study,s6review}.}
{fig:towerdiag}

\subsection{Reading the original breathing row}
The original \(\epsilon=1\) row records final mean \(|R|=0.5489832686504329\), four surviving bands, mean \(L=11.905260171214326\), population std \(0.9892193445854098\), range \(3.5626358780586127\), and dominant frequency \(1/150\). All seven rows share final mean amplitude and band count. Summary equality alone would not prove path equality.

Retained R1 gives stronger path-specific evidence: selected \(\epsilon=0/1\) pairs have identical saved \(X,V,R\) paths and no clipping activation. Its dominant bin follows changed window lengths. The historical driver uses \(L[50:]\): 150 readouts at passes 51--200, reading states 50--199. Its detrending subtracts only a mean. After a std guard it flags oscillation if the largest non-DC power exceeds 1\% of total. In exact mean-zero arithmetic 75 non-DC bins have a maximum at least \(1/75>1\%\) of nonzero total. This reviewed argument identifies a weak oscillator test, not a new FFT result.

The affine readout, inactive clipping on tested paths, and window-dependent bin do not establish an independent self-sustaining 150-pass field. Earlier R2 controls remain retained context for additions and compaction; none is rerun. Separate R3 material retains its review status and is not promoted to accepted science.
\section{Universe: source-faithful construction}\label{sec:universe}
\subsection{Initialization retains discarded random draws}
Universe's constructor (\src{U}{60--96}) stores \(X,R,V,p\), configuration and fixed frequencies, but no Tower II band ID. With multiple frequencies enabled it independently assigns
\begin{equation}\label{eq:Ufreq}
\omega_i\in\{0.244,\;0.244\phi,\;0.244\phi^2,\;0.244/\phi\}.
\end{equation}
With the option off, every label is 0.244 and no choice draw occurs. Allocation is not balanced. Stored \texttt{gamma\_srg=0.08699} is unused throughout this file. Despite its Infinity Part II header, this is a one-real-amplitude engine, not continuation of Infinity's two-field state.

The selected driver (\src{U}{675--702}) seeds with 1618 and draws six 20-point normal clouds, coordinate standard deviation 1.5, stacked in center order
\[
(8,0,0),\;(-8,0,0),\;(0,8,0),\;(0,-8,0),\;(0,0,8),\;(0,0,-8).
\]
It calls the constructor with \(N=120,h=6\), multiple frequencies and breathing enabled, \texttt{seed=None}. The constructor still consumes uniform \(X\), uniform \(R\), then frequency-choice draws. The driver replaces \(X\) by the clouds and consumes a fresh uniform amplitude draw. Discarded constructor arrays remain part of stream history; omitting them changes later draws.

The accepted replay preserved this order. Counts in sorted frequency order are 29,36,28,27, and each population happens to contain all four frequencies. This is a realized allocation, not a rule. Population \(P_k\) means IDs \(20k,\ldots,20k+19\); the recurrence never consults that label.

\subsection{Incoming selection and synchronous amplitude stage}
Universe computes incoming distances once (\src{U}{105--125}). Finite \(h\) filters strict \(d_{ij}<h\), explicitly removes self, sorts eligible distances and chooses at most five. With \texttt{None}, all non-self rows are eligible and self is still explicitly removed. Finite rows may be short or empty. An unrestricted minimum-six component argument cannot be transferred to arbitrary horizon-filtered graphs.

Local movement is the simultaneous Equation~\eqref{eq:local}. The amplitude helper still reads incoming \(R\) and distances, although positions have become \(Y\). A nonempty row first normalizes inverse distance:
\begin{equation}\label{eq:Uweights}
u_{ij}=(d_{ij}+10^{-8})^{-1},\qquad
\bar u_{ij}=u_{ij}/\sum_{k\in\mathcal N_i}u_{ik}.
\end{equation}
With multiple frequencies, form
\begin{equation}\label{eq:Uguard}
v_{ij}=\bar u_{ij}e^{-5(\omega_j-\omega_i)^2},\quad
s_i=\sum_jv_{ij},\qquad W_{ij}=v_{ij}/s_i
\quad\text{only if }s_i>10^{-10}.
\end{equation}
Otherwise skip the row. Equality skips here, unlike Tower II's raw-sum guard. In single-frequency mode \(W=\bar u\) and the second guard is absent. There is no additional guard on the first normalization.

For active rows,
\begin{equation}\label{eq:UB}
B_i=\left[R_i+0.15\left(\sum_jW_{ij}R_j-R_i\right)\right]
\begin{cases}
1+0.01\sin(2\pi\omega_i p),&\text{multiple frequencies},\\
1,&\text{single frequency}.
\end{cases}
\end{equation}
Skipped rows retain incoming \(R_i\) and skip sine. Every row reads old \(R\); installation is synchronous. There are no Tower II phi-target or sum/difference additions. Golden spacing enters the catalogue, but the kernel uses ordinary squared frequency difference rather than log ratio. Equation~\eqref{eq:Q} then supplies \(Q\).

\subsection{Ascending local fusion: which coordinates are read}
Let
\[
a_{ij}=\frac{|Q_j|+10^{-8}}{\sum_{k\in\mathcal N_i}(|Q_k|+10^{-8})},
\qquad \kappa_i=0.03Q_i.
\]
Let \(Z_j\) denote a row's completed fusion write. For finite horizon (\src{U}{178--188}), ascending \(i=0,\ldots,N-1\) gives, for nonempty rows,
\begin{align}
c_i&=\sum_{\substack{j\in\mathcal N_i\\j<i}}a_{ij}Z_j
     +\sum_{\substack{j\in\mathcal N_i\\j>i}}a_{ij}Y_j,\label{eq:orderedcenter}\\
Z_i&=Y_i+\kappa_i(c_i-Y_i).\label{eq:orderedwrite}
\end{align}
Empty rows give \(Z_i=Y_i\). Direct self is excluded. Advanced indexing creates a fresh selected-coordinate operand for each center; it is not a persistent view. Later centers select again from the then-current position array.

Earlier IDs are read after their writes, later IDs before them. \(Q\) and neighbor identities remain fixed through fusion. Centers use magnitude-amplitude weights with no distance factor; displacement uses signed \(Q_i\). A row's own coordinate is still \(Y_i\), because prior rows wrote only themselves.

An earlier \(Z_j\) contains \(\kappa_j(c_j-Y_j)\); substitution exposes indirect selected paths within a pass. Indirect self-dependence is possible when \(i\) was selected by an earlier row, despite direct self-exclusion. This is the reviewed dependency expression, not a new convergence or general order-sensitivity theorem. Coefficients or cancellations can make a contribution vanish.

The horizon restricts pairs at the incoming boundary. Distances are not re-tested after local movement or earlier writes. Selected chains can span farther than \(h\); there is no one-hop-per-pass physical speed law. For fixed selections, disconnected weak components have no additional all-point or band-center bridge in this branch. This conditional structural fact does not establish permanent isolation as geometry changes.

\subsection{Global branch, prescribed bound and final record}
With horizon \texttt{None}, one all-point center, including self, is computed before a vectorized displacement:
\begin{equation}\label{eq:global}
c=\frac{\sum_j(|Q_j|+10^{-8})Y_j}{\sum_j(|Q_j|+10^{-8})},
\qquad Z_i=Y_i+0.03Q_i(c-Y_i).
\end{equation}
Thus switching to \texttt{None} changes fusion as well as eligibility. It does not make local amplitude exchange all-to-all.

After fusion,
\begin{equation}\label{eq:LU}
L_p=\begin{cases}
9[1+0.3\sin(2\pi\,0.244p)e^{-0.01p}],&\text{breathing enabled},\\
9,&\text{otherwise}.
\end{cases}
\end{equation}
Compression, noise and \(V\) follow Equations~\eqref{eq:rpco}--\eqref{eq:noise}. \(L_p\) is time-only, unlike Tower II's incoming-amplitude readout. Neither bound is a coordinate wall. Observers (\src{U}{205--259}) receive old adjacency/distances but do not use them; they measure completed arrays.

\section{Universe's original outputs}\label{sec:original}
The five historical tables have explicit source writes, matching schemas and driver settings. They contain 233 data rows, 3,307 cells: 3,291 finite numeric cells, 15 text cells and one missing decay-time value. Original strings remain unchanged~\citep{s7cstudy,s7creview}.

{\small
\begin{longtable}{P{40mm}P{28mm}P{78mm}}
\caption{\ev{ORIGINAL / ATTRIBUTED}. All five outputs and generating drivers. Rows exclude headers; physical CSV line = row \(+1\).}\label{tab:drivers}\\
\toprule Table & Driver / write & Configuration and meaning\\\midrule
\endfirsthead\toprule Table & Driver / write & Configuration and meaning\\\midrule\endhead
\protect\path{full_universe.csv} & U:663--758 / 757 & \(200\times15\). Seed 1618 selected setting; completed passes 1--200. No initial row or individual coordinates.\\
\protect\path{causal_horizons.csv} & U:378--497 / 484 & \(18\times11\). Horizons 3,5,8,12,20,None, each seeds 42--44. Four normal populations, \(N=80\), 100-pass endpoints; frequency/breathing off.\\
\protect\path{condensate_perturbation.csv} & U:273--371 / 358 & \(7\times7\). Seven kicks after 100-pass \(N=50\), seed-2038 preparation; 50 response passes each. Global fusion, single frequency, static bound.\\
\protect\path{multi_frequency.csv} & U:504--576 / 575 & \(6\times7\). Single/multiple frequency, seeds 42--44, \(N=60\), 150 passes. Global fusion, static bound.\\
\protect\path{breathing_condensate.csv} & U:583--656 / 654 & \(2\times9\). Static/breathing, \(N=50\), seed 2038 reset per setting, 200 passes. Global fusion, single frequency.\\
\bottomrule
\end{longtable}}

\subsection{Main historical chronology}
The spatial observer reports 49 groups at pass 1, 12 at 9, 11 at 14, and ten at 16; ten persists through 200. Before replay, this established stable count, not membership. Amplitudes remain unequal: at pass 20 the amplitude-gap observer gives one group despite max/min \(3.060341\). At 200, mean \(|R|=0.144070880527\), signed std \(0.040448472042\), ratio 0.280754. The ratio reached about 0.252569 at pass 13, so progressive exact equalization is not the recorded history.

Mean GRF stays below one, ending at 0.994912594. Its definition supports some nonzero displacement each pass, not individual directions or recovered mean speed. Four frequency columns are whole-field means of supplied labels, not spatial groups. Figure~\ref{fig:orig} is the accepted later rendering of these original aggregates.

\plate{evidence/S7C/source_package/figures/F01_original_universe_chronology.png}
{\ev{OBSERVED IN HISTORICAL DATA / DERIVED}. All original main-table rows 1--200; no coordinates inferred. Panels distinguish amplitude extrema, assigned-frequency means, two grouping observers, mean reciprocal displacement, prescribed \(L_p\), and the accepted passive \(R_{\rm std}/R_{\rm mean}\) ratio. The six-population reference comes from the driver. “New” in the unchanged image refers to S7C's earlier rendering, not a present calculation. S7C F01~\citep{s7cstudy}.}
{fig:orig}

\subsection{Horizon settings are not time steps}
The driver places 20-point normal populations with coordinate std 1 near \((5,5,5)\), \((5,-5,-5)\), \((-5,5,-5)\) and \((-5,-5,5)\). Each trial resets its seed inside each horizon setting, retains discarded constructor draws, replaces \(X,R\), then evolves 100 passes.

All 15 finite-horizon rows report five spatial groups, while the three \texttt{None} rows report four. Mean pair distances between the four original-population arithmetic centers lie about 13.87555--14.25550 for finite settings and 8.22412--9.17227 for \texttt{None}. Within a trial, stored values other than horizon agree for \(h=5,8,12,20\). These are historical summary findings, not full-state or graph identity. The local/global fusion switch prevents a radius-only causal comparison.

Sentinel 999 means \texttt{None}, not radius 999. The flattened-coordinate intra-spread statistic is distinguished from radius in Section~\ref{sec:pooled}.

\plate{evidence/S7C/source_package/figures/F02_horizon_settings.png}
{\ev{OBSERVED IN HISTORICAL DATA}. All 18 horizon/trial endpoints after 100 passes. The x-axis is categorical; offsets separate trials. No line represents time between settings. The division before None marks the switch from ordered local fusion to one global center. Distances are between original-population arithmetic centers in uncalibrated coordinates. S7C F02, unchanged~\citep{s7cstudy}.}
{fig:horizons}

\subsection{Perturbation, frequency and breathing summaries}
The perturbation driver saves \(X,R\) after 100 passes, then restores those arrays, zeros \(V\), clears history and resets \(p=100\) before each kick to ID 0. It does not restore RNG state: successive strengths use successive noise segments. Initial deviation means the first recorded post-kick state at 101, not the kick itself; final deviation is absolute deviation at 150.

All seven rows are \texttt{DECAYING}, with final deviation smaller than first recorded deviation. The 0.01 kick has one absolute-deviation peak above 0.001; other rows have zero. The 0.05 row leaves decay time blank because the sampled 10\% threshold was never reached. The ordered classifier tests final \(<0.001\) times maximum recorded deviation (ABSORBED), more than three peaks (OSCILLATING), final \(>0.5\) times first deviation (PERSISTENT), else DECAYING. Temporary response histories were not written, so these classifier summaries cannot now be independently reconstructed. The narrative's all-size absorption and excitation-spectrum interpretation exceeds the seven finite responses.

The multi-frequency table reports four supplied labels and final band-mean ratios \(1.10766886\), \(1.27975898\) and \(1.05919374\). Label count is not a survival test or a newly generated spectrum. Toggling the option changes RNG consumption and enables sine as well as the kernel.

Static/breathing rows have identical exported numbers, including \(R_{\max}=0.179158800380\), signed std \(0.019352317340\), and four amplitude-gap groups. This is equality of summaries, not all states or cause. This driver fits a quadratic to all 200 maxima, then takes an rFFT of residual rows 51--200, unlike Tower II's mean-only subtraction. Saved \(1/150\) and fraction 0.91888455 remain historical calculations, not a new oscillator diagnosis. Main replay clipping conclusions are not transferred to this distinct driver.

\section{Replay fidelity and historical agreement}\label{sec:fidelity}
The accepted reference extracted helpers, complete class and selected driver from U:42--54,60--266,663--758, with necessary imports and run-local output path. Other drivers and original main were not executed. Capture used the same operations plus copy-only hooks; removing them recovered the reference text and definitions. Original bytes remained unchanged~\citep{s7rprovenance,s7rreview}.

The earlier environment used NumPy 2.4.4, pandas 3.0.3, SciPy 1.17.1 in \texttt{torment}. Records include every boundary, incoming selection, amplitude stage, center operand, row write, actual random return and observer. Each run has 211 random-interface calls: one seed, ten value-producing initialization calls, 200 noise returns. This counts API calls, not generator words. Before-seed arbitrary process state is excluded from equality; deterministic checkpoints from seed onward agree.

Reference and capture match in dtype, shape, raw array bytes, original metric bits, random returns and deterministic RNG checkpoints. Their exports are byte-identical. This verifies instrumentation in one setting, not independent trials or cross-platform replication. Accepted independent reviews checked fixed-record stages without propagating a new trajectory; these checks were not rerun for this paper.

Historical comparison is separate. The two main tables have 3,000 cells including 200 pass labels, hence 2,800 non-pass metric cells. All 800 integer cells match exactly. All 2,200 continuous cells satisfy
\begin{equation}\label{eq:tolerance}
|r-h|\le10^{-12}+10^{-10}|h|,
\end{equation}
where \(h\) is the historical value. Exactly 2,312 parsed cells agree: 800 integer plus 1,512 continuous. The remaining 688 differ slightly; maximum absolute discrepancy is \(1.7763568394\times10^{-15}\), in entropy. Historical/replay CSV bytes differ. Agreement does not identify later coordinates with an unavailable original path.

\plate{evidence/S7R/source_package/figures/F05_metrics_and_amplitudes.png}
{\ev{REPRODUCED}, compared with \ev{ORIGINAL} aggregates. S7R F05, unchanged, passes 1--200. Upper left is maximum historical discrepancy divided by frozen tolerance; acceptance level one is far above plotted ratios. Other panels show positive signed amplitudes, fixed-frequency means and prescribed bound. Zero clipping is a captured pointwise result, not inferred from the bound's appearance~\citep{s7rprovenance,s7rreview}.}
{fig:fidelity}

\section{Universe motion, groups and actual reads}\label{sec:umotion}
\subsection{Concentration within separated populations}
Whole \(R_g\) changes from \(8.316796374\) to \(8.071505534\), diameter from \(22.398761997\) to \(19.021944473\). Large spread remains because populations remain separated. At boundary 16 \(R_g\approx8.049319139\), less than the endpoint; endpoint reduction is not monotonicity.

Population contraction differs: P3 at 200 has \(R_g\approx0.0669491\), diameter 0.166309; P5 has \(R_g\approx1.80620\). P5 contains two stable observer groups, so its radius includes separation between pieces. P4's arithmetic center translates about 0.898055 from initialization, P3's about 0.305873. Center translation, spread and grouping are distinct~\citep{s7rmotion,s7rreview}.

\plate{evidence/S7R/source_package/figures/F01_world_snapshots.png}
{\ev{REPRODUCED}. Universe seed 1618, boundaries 0,16,50,200, 120 persistent IDs. Same camera, equal aspect, \(\pm13\) world axes. Color/shape denotes supplied populations; labels identify selected stable rows, including 97. Overlapping markers retain different identities. S7R F01, unchanged; no original coordinate movie is implied~\citep{s7rmotion}.}
{fig:uworld}

\subsection{Stable membership does not freeze the graph}
Ten observer member sets remain unchanged over 185 observations, passes 16--200, with sizes \(10,10,20,9,11,20,9,11,13,7\):
\begin{center}\small
\begin{tabular}{lrrrrrr}\toprule
Population & P0 & P1 & P2 & P3 & P4 & P5\\\midrule
Group sizes & \(10+10\) & 20 & \(9+11\) & 20 & \(9+11\) & \(13+7\)\\
\bottomrule\end{tabular}
\end{center}
No late observer group mixes supplied populations. This adds membership information to historical count information, but remains reproduced evidence.

At the \emph{same boundary 199}, observer/weak/strong counts are ten/13/14. P1's one 20-point observer group contains weak pieces \(8+6+6\); P3's contains \(12+8\). At boundary 16 the opposite mismatch also occurs: P0 has two ten-point observer groups but one weak local component of 20. At that same boundary P1 already has one observer group and three weak components. Neither partition universally refines the other.

All 200 captured incoming graphs have 600 edges and no empty rows in this realization. Across 199 successive graph comparisons, 35,682 binary entries change; the final comparison, boundaries 198 to 199, changes 244. These are changes counted across time, not unique pairs or independent experiments. Stable memberships coexist with rearranged selections.

\wideplate{evidence/S7R/source_package/figures/F02_group_memberships.png}
{\ev{REPRODUCED / DERIVED}. S7R F02 compares supplied populations, frequencies, weak/strong components and spatial groups at boundaries 16 and 199. Graph at \(b\) supplies \(b+1\); observer measures \(b\). Colors are categorical within rows, not identity across definitions. P0/P1 at 16 demonstrate mismatch in both directions. No final-boundary-200 graph is claimed~\citep{s7rreview}.}
{fig:partitions}

\subsection{Brief contact and a turning path}
Only three directed pairs cross populations: 39 selects 97 at passes 1--3, 97 selects 34 at 1--6, and 97 selects 39 at 1--8. These total 17 visits between P1 and P4. None occurs in input graphs supplying 9--200. U:693's comment that each cluster only sees itself is false for these early replay selections.

GPT's additional readback separates leaving a radius from losing a nearest-five place:
\begin{table}[H]\centering\small
\begin{tabular}{lrrr}\toprule
Selection & First absent pass & Pair distance & Farthest selected\\\midrule
39 selects 97 & 4 & 3.020634879 & 2.810868773\\
97 selects 34 & 7 & 4.836983324 & 4.808048733\\
97 selects 39 & 9 & 4.840333657 & 4.628762859\\
\bottomrule\end{tabular}
\caption{\ev{REPRODUCED / DERIVED}. S7R \protect\path{independent/FINDINGS.json}, \texttt{cross\_selection\_departure}. All departing pairs remain inside \(h=6\); five closer candidates occupy the selected places.}\label{tab:depart}
\end{table}
This accounts for actual selection loss, not the trajectory effect of forcing a link to remain.

ID 97 starts near \((-5.97170,-0.946781,5.79831)\). Its \(z\) first falls then rises while \(x\) moves toward P4: boundary 8 is near \((-4.15099,-1.16281,5.13300)\), 16 near \((-1.78667,-1.36076,6.13502)\), and 200 near \((0.394820,-1.52310,6.37897)\). Its sampled path length \(8.119170932\) is the largest among 120 IDs. It ends with IDs 83,85,86,88,91,92,94,95,96,97,99, all P4. This is a sampled turn, not wrapping or identity exchange; no explorer role follows.

\wideplate{evidence/S7R/source_package/figures/F06_transient_cross_population_contact.png}
{\ev{REPRODUCED / DERIVED}. S7R F06, unchanged. Left: sampled \(x\)--\(z\) paths of 34,39,97 over 0--30, omitting \(y\). Middle: incoming distances with actual-selection markers; radius alone does not determine selection. Right: ID 97's completed steps through 200. Joins are reading guides, not a continuous-time law. Table~\ref{tab:depart} provides first-absence distances~\citep{s7rmotion,s7rreview}.}
{fig:cross}

\subsection{Earlier writes enter later centers}
The 24,000 fusion row visits contain 120,000 selected-coordinate reads: 60,118 earlier-written and 59,882 later-unwritten. Weighted earlier-write contribution is nonzero in 19,974 centers. Counts cover one path, not independent experiments.

Let \(X^{\rm read}_{ij}\) be the actual coordinate passed for neighbor \(j\). The reviewed accounting identity is
\begin{align}
d_i^{\rm read}&=\sum_{j\in\mathcal N_i}a_{ij}(X^{\rm read}_{ij}-Y_j),\label{eq:dread}\\
c_i&=\sum_{j\in\mathcal N_i}a_{ij}Y_j+d_i^{\rm read}.\label{eq:decompose}
\end{align}
It decomposes one actual center; floating reduction differences leave small reported residuals. It is not a simultaneously evolved alternative, nor an endpoint difference between algorithms.

Accepted representative rows maximize \(\norm{d_i^{\rm read}}\) within each fixed pass 1,16,200:
\begin{center}\small
\begin{tabular}{rrlr}\toprule
Pass & Row & Earlier-written selected IDs & \(\norm{d_i^{\rm read}}\)\\\midrule
1 & 112 & 100,102,103,107,109 & .008262981568\\
16 & 6 & 2,5 & .001045505511\\
200 & 8 & 0,1,3,7 & .0000213994194\\
\bottomrule\end{tabular}
\end{center}
This selected recorded operands, not seeds or parameters.

At pass 16 row 6 selects \(2,5,11,16,18\). IDs 2 and 5 have been written; the others are still \(Y\). In that order, normalized weights are approximately 0.2408018130, 0.2448881163, 0.1692338522, 0.1773596550 and 0.1677165636. For ID 5,
\begin{align*}
Y_5&=(10.1642402234,-0.009263063370,-1.0188735915),\\
X^{\rm read}_{6,5}&=(10.1645760721,-0.008571991931,-1.0211011336).
\end{align*}
The total earlier-write contribution is
\[
d_6^{\rm read}\approx(0.000063805420,\ 0.000308562427,\ -0.000996895115).
\]
The actual center is approximately \((10.2112525079,0.123574916694,-1.4044326301)\). With \(Q_6=0.145515333383\), row 6's fusion displacement is approximately
\[
(-0.000157993719,\ 0.000573960922,\ -0.000788661172).
\]
These are pre-noise, not completed-state quantities. Maximum accounting-residual component is \(1.77636\times10^{-15}\). Exact operands are in the stage NPZ and readable episode JSON~\citep{s7rordered,s7rreview}.

\wideplate{evidence/S7R/source_package/figures/F04_actual_ordered_reads.png}
{\ev{REPRODUCED / DERIVED}. S7R F04: pass 1 row 112, pass 16 row 6, pass 200 row 8. Upper sampled paths cover 0--6,11--21,195--200; each has one fixed translation, equal axes and a disclosed differing close-view scale. Circles start, triangles end. Lower bars are signed weighted actual-read contributions; stars mark earlier writes. Later rows contribute zero to Equation~\eqref{eq:dread}. These are not reordered controls~\citep{s7rordered}.}
{fig:ordered}

\subsection{Continuing motion and stage scales}
All 120 points move at pass 200: mean step \(0.005118790595\), minimum about \(0.000827802018\), maximum \(0.011059750237\). Mean cumulative sampled path is \(3.481760184\). Stable membership is compatible with these small motions.

At pass 1 mean lengths of local, fusion and noise increments are \(0.203511,0.00562931,0.0168112\); net mean step is 0.209809. At 200 the stage means are \(0.00156648,0.0000274445,0.00496697\), net 0.00511879. The largest endpoint stage mean is recorded noise. This describes operands, not a noise-off experiment. Vector increments add; scalar lengths generally do not.

\plate{evidence/S7R/source_package/figures/F03_geometry_and_movement.png}
{\ev{REPRODUCED / DERIVED}. S7R F03 separates centered population spread, center translation, stage lengths and grouping definitions. Geometry: boundaries 0--200; increments: passes 1--200; graphs: incoming 0--199; observers: completed 1--200. Named member sets are measurements, not rigid bodies. Small late motion persists despite stable spatial membership~\citep{s7rmotion}.}
{fig:umotion}


\section{Observers, bounds and reviewed mathematics}\label{sec:observers}
\subsection{What the labels measure}
Universe sorts \(|R|\) descending, counts adjacent gaps whose magnitude exceeds \(0.05\max|R|\), and adds one. One amplitude group means no such large adjacent gap; it does not require equal extrema. Spatial groups come from average linkage of completed \(X\), cut at distance 2.0. Neither observer drives the next update.

With \(\pi_i=|R_i|/\sum_j|R_j|\), entropy is
\begin{equation}
\mathcal S=-\sum_i\pi_i\log(\pi_i+10^{-12}),
\end{equation}
or zero if the amplitude sum is zero. This is a dimensionless distribution statistic. Maximum, minimum and mean use \(|R|\), but \texttt{R\_std} uses signed \(R\). \texttt{top\_seed\_id} is the first argmax of \(|R|\), not a leader trajectory. Four-decimal frequency keys label fixed assignments, not measured instantaneous frequency.

Both sources' GRF observer is
\begin{equation}\label{eq:grf}
\overline{\mathrm{GRF}}=\frac1N\sum_i\frac1{1+\norm{V_i}},
\end{equation}
not \(1/(1+\mean_i\norm{V_i})\). This relation to displacement defines no force law. Lock, cluster, separate, merged and condensate must be read through actual thresholds. The full-universe printout calls population centers separate when distance exceeds 3; its diversity threshold is not a proof of independence.

\subsection{A large pooled spread can describe a tiny cloud}\label{sec:pooled}
The horizon and full-universe drivers use \texttt{np.std(pos\_k)} with no axis. For \(n\) points with vector mean \(\mu=(\mu_x,\mu_y,\mu_z)\), this pools \(3n\) entries around scalar mean \(m=(\mu_x+\mu_y+\mu_z)/3\). The reviewed identity is
\begin{equation}\label{eq:pooled}
s_{\rm flat}^2=\frac1{3n}\sum_{i,d}(X_{id}-m)^2
=\frac{R_g^2}{3}+\frac13\sum_d(\mu_d-m)^2.
\end{equation}
Write \(X_{id}-m=(X_{id}-\mu_d)+(\mu_d-m)\). Summing deviations over \(i\) gives zero, eliminating the cross term. The squared first term gives \(R_g^2/3\), the repeated center term the remainder. This expands the accepted proof; no new statistic is calculated.

A translation can change differences among the center's coordinate components without changing \(R_g\). Pooled spread need not be invariant under arbitrary translation/rotation or measure cloud size. For reproduced P3 at 200 it is \(3.238809712\), while centered \(R_g=0.066949097\) and diameter \(0.166308988\). Both definitions and numbers are retained. The centered statistic is not substituted into an old table and called original evidence.

\subsection{Inactive clipping remains a scoped observation}
Equation~\eqref{eq:rpco} sends zero to zero. For nonzero \(q\) inside its logarithmic bounds, exact real arithmetic gives
\begin{equation}\label{eq:eps}
\operatorname{RPCO}_L(q)=q+10^{-9}\sgn(q).
\end{equation}
Inactive clipping does not remove epsilon or floating \(\log/\exp\). The lower floor can raise small nonzero magnitudes and the upper limit can cap large ones.

Tower II's selected conclusion is supported by retained R1 paired controls. Universe's is supported by actual flags in one replay: zero lower/upper activations among 24,000 point updates, minimum log-space margin about \(3.032068826\). All saved signed \(R\) and captured \(Q\) were positive. No Universe breathing-off control was run. This identifies the branch used on that path, not universal irrelevance or the cause of another driver's equal summaries.

The identities and amplitude diagnostics introduce no energy minimum, variational law or conservation claim. They make measurements and dependencies precise enough to interpret the source terminology.

\section{Comparison and historical meaning}\label{sec:comparison}
Shared local movement, \(H\), signed compression and declining noise do not make the engines identical. Their centers and amplitude rules create different dependencies.

\begin{longtable}{P{25mm}P{60mm}P{62mm}}
\caption{\ev{DERIVED FROM ORIGINAL SOURCE}. Selected mechanisms, anchored to T:62--254 and U:60--266.}\label{tab:compare}\\
\toprule Mechanism & Tower II & Universe\\\midrule
\endfirsthead\toprule Mechanism & Tower II & Universe\\\midrule\endhead
State & \(X,R,p\), fixed band/frequency; diagnostic \(V\) & \(X,R,p\), fixed frequencies; diagnostic \(V\); no band ID or second amplitude\\
Initialization & Uniform box; \(N=80\), seed 2038; balanced shuffled bands & Six Gaussian populations; \(N=120\), seed 1618; discarded constructor draws retained; independent labels\\
Selection & Unrestricted sort/slice assumes self first; finite branch explicitly filters self & Strict \(h=6\), explicit non-self, nearest-five cap; short/empty rows possible in general\\
Local movement & Simultaneous \(0.05\ell+0.15\tanh\ell\) & Same simultaneous stage\\
Amplitude weights & Inverse distance times log-ratio preference; raw-sum guard skips \(<10^{-10}\) & Normalize inverse distance; squared frequency-difference factor; second guard requires \(>10^{-10}\)\\
Sine / additions & Assigned sine, then additive phi-target and sum/difference terms & Assigned sine when enabled; no additive mixing\\
Center & All fixed-band members, self included; once per band & Selected neighbors, directly excludes self; recomputed per ascending ID from current coordinates\\
Coefficient & Signed \(0.02Q_i\) & Signed \(0.03Q_i\)\\
Global option & Selected construction has no single all-point center & None switches to all-point, once-computed center and simultaneous fusion\\
Bound & Incoming-amplitude readout \(9(1+\epsilon\mean|R|)\) & Prescribed \(9[1+0.3\sin(2\pi0.244p)e^{-0.01p}]\)\\
Clock / retention & Integer pass; \(V\) not consumed; no wrapping & Same; no velocity damping or inertia\\
Observers & Fixed-band coherence/first lock, amplitude summaries and band GRF & Amplitude gaps, average linkage, entropy, frequency means and GRF\\
Evidence & Original summary at 200; R1 0--600, extension after 200 & Original aggregates 1--200; accepted twins 0--200\\
\bottomrule
\end{longtable}

Tower II's strong concentration and Universe's retained large-scale separation occurred under different geometry, \(N\), seeds, labels, kernels, additive terms, schedules, centers and durations. They are not a controlled one-factor experiment. Source establishes different algorithms; records establish behavior under specified settings. Contraction differences cannot be attributed solely to band-wide versus local fusion.

Accepted SIM04/SIM05 context describes selected engines using one all-point center; Infinity also carries a second real amplitude and cross terms. This locates the present mechanisms without reopening Volume II. Universe's historical cross-reference to Infinity does not make it the same evolution. Shared constants and pictures do not establish one canonical SRG model.

Tower II's note asks whether lock predictions, self-coupled breathing, frequency content and tunneling reveal structure. Universe's note asks how perturbations, horizons, assigned frequencies and breathing enrich a seemingly simple condensate. These are attributed motivations, preserved respectfully. The reconstruction distinguishes questions asked from outcomes established: inserted catalogues explain label persistence without proving emergence; finite amplitude-gap counts do not establish a uniform infinite-time ground state.

\section{What this volume closes and leaves open}\label{sec:limits}
Tower II is recoverable stage by stage, with its selected fixed-record motion tied to persistent IDs, radii, local components and band dependencies. Universe's source-to-output map is recoverable, and its accepted finite replay supplies actual coordinates and within-pass operands beyond original aggregates. The three reviewed identities explain band variance, ordered reads and pooled coordinate spread without replacing historical measurements.

Original individual trajectories and historical runtime environments remain unrecovered. Finite reproduced paths do not prove infinite-time convergence or permanent independence. Actual ordered reads do not measure a reversed/simultaneous control. Other settings' clipping and the necessity of individual stages are not determined here. Separate R3 material retains review status.

SIM06 and SIM07 are complete within accepted scopes. Published Volumes I/II remain frozen. This candidate does not reconstruct every historical SRG family. The \textbf{BreathingField/field-merger study is deferred and outside this volume}. Drafting now supersedes earlier sequencing without editing old plans. The next step is GPT and Hilmir's review, not an automatic experiment or publication.

\appendix
\section{Source and evidence navigation}\label{app:evidence}
Paths are relative to the review package. T is
\protect\path{evidence/S6/source_package/original/srg_golden_tower_II.py};
U is
\protect\path{evidence/S7C/source_package/original/srg_condensate_breaking.py}.
They preserve original line numbering. T's SHA-256 is
\protect\path{147a32be3b6d6aad17ada6ab7be2fbc3994c3363f74d6a2b351fb8eecda80304};
U's is
\protect\path{a645a3bb564a2cae5e2b3ac3744f3f34f559a67132f255b3b0a97d277e68beb5}.

\begin{longtable}{P{31mm}P{49mm}P{68mm}}
\caption{Exact source locator map.}\\
\toprule Object & Tower II lines & Universe lines\\\midrule
\endfirsthead\toprule Object & Tower II lines & Universe lines\\\midrule\endhead
Helpers, constants & 31--46 & 42--54\\
Constructor & 62--109 & 61--96\\
Neighborhood & 111--127 & 105--125\\
Bound & 135--138 & 98--103; used 195\\
Local movement & 140--148 & 160--168\\
Amplitude stage & 150--190 & 127--153; installed 171\\
Common \(H\) & 192--194 & 173--176\\
Fusion & 196--206 & 178--192\\
Compression, noise, \(V\) & 209--216 & 194--203\\
Observers/evolve & 218--254 & 205--266\\
Selected driver & 376--472 & 663--758\\
Other drivers & 262--368;480--821 & 273--656\\
Main & 828--842 & 765--778\\
\bottomrule
\end{longtable}

Three complete accepted ZIPs are in \protect\path{handoff/inputs/}. The accompanying evidence index provides relative links, archive hashes and exact member locators. Extract each archive into a separate directory to follow inherited report links; never flatten their source-package roots.

In S6, arrays are under \protect\path{source_package/records/R1_seed2038_eps1/}.
\texttt{X.npy}, \texttt{R.npy}, \texttt{V.npy} have boundaries 0--600;
\texttt{Q.npy}, \texttt{L\_used.npy}, \texttt{positions\_at\_rpc.npy},
\texttt{actual\_rpc\_output.npy}, \texttt{noise.npy} index passes 1--600.
The accepted union table is \protect\path{independent/combined_dependency_graph.csv};
geometry and selected neighbors are under \protect\path{source_package/derived/}.
No new scientific calculation on these arrays was performed.

In S7R, \protect\path{source_package/runs/A_REFERENCE/} and
\protect\path{source_package/runs/B_CAPTURE/} each contain 201 boundary NPZs.
Capture \protect\path{stages/pass_016.npz} retains row-6 operands; readable vectors are
\protect\path{source_package/analysis/ORDERED_EPISODES.json}.
GPT's counts are \protect\path{independent/GRAPH_COUNTS.csv} and
\protect\path{independent/ORDERED_READ_COUNTS.csv}.
The full schema is \protect\path{evidence/S7R/source_package/RECORD_SCHEMA.md}.
Binary boundaries, actual random returns and stages remain in the archive without a second extracted trajectory tree.

\protect\path{CLAIM_EVIDENCE_MATRIX.md} maps claims to originals, accepted reports, data and equations.
\protect\path{FIGURE_REGISTER.json} records unchanged assets and original plotting specifications.
Printed rounding never overwrites stored values.

\section{Review incorporation and document reproducibility}\label{app:review}
S6 supplies the late band-0 bridging example and all-600-boundary weak connectivity of the union, retaining the R1 extension and R3-status limits. S7C supplies driver attribution, observer/geometry distinctions and pooled-variance explanation. S7R supplies strict fidelity, cell accounting, same-boundary partition mismatch in both directions, nearest-five link loss within the horizon, and earlier/later read totals. All are retained accepted results; no scientific verification was rerun.

Current work consists of source/document inspection, file identity checks, unchanged evidence copies, typesetting and document QA. The final page-QA record identifies every rendered page. Build instructions invoke document operations only. Archived simulators, scientific methods and old work orders are evidence, not new authority to execute.

One process deviation is disclosed: initial attachment and instruction-file reads used PowerShell before the CMD-only provision was read. Subsequent document work used Windows CMD, with the existing torment environment for Python. No scientific operation occurred during initial reads. This does not meet the requested CMD-throughout procedure and is recorded instead of concealed or treated as reason for a scientific rerun.

Preservation rechecks the actual attachment, accepted nested archives and copied evidence against recorded identities. It does not assert a fresh whole-repository audit. Historical sources, accepted studies, canonical ledger, kernels, unrelated work and published editions were not edited. Payload manifest and external ZIP identity establish review-file integrity; hashes alone do not prove non-execution.


\clearpage{\raggedright\bibliographystyle{plainnat}\bibliography{references}}
\end{document}

